\chapter{Introduction}
\label{ch:introduction}
% Port from: docs/01_overview.md

\section{Background}
Nepal's cultural landscape is dense with sites worth exploring: the country
holds ten UNESCO World Heritage properties, four of them cultural sites within
the Kathmandu Valley alone, and it welcomed 1.147 million international
tourists in 2024. Yet the experience of actually visiting such a site --- or,
at a smaller scale, of a newcomer finding their way around a university
campus --- remains largely unassisted. Navigation applications reliably bring
a visitor to the entrance of a monument or an institution, but the guidance
ends at the gate. Inside, visitors depend on scarce physical tourism centres,
sparse signage, or the availability of a human guide.

At the same time, the technical preconditions for a software answer to this
gap have matured. Modern smartphones ship with GPS receivers, magnetometers,
and cameras as standard hardware, and web browsers now expose geolocation,
device orientation, and camera access directly to web pages, making
location-based Augmented Reality (AR) possible without installing a native
application. \projectname{} was built on this premise for Orchid HackX 2026
by Team Valhalla: a guided, gamified AR walk that begins where conventional
navigation stops.

\section{Project Overview}
\projectname{} is a mobile-first, gamified AR navigation platform for
educational campuses and heritage sites. A user opens the web application in
a smartphone browser, selects a pre-mapped route, and is guided in real time
by a live compass arrow together with a GPS distance readout to the next
waypoint. When the user physically arrives at a checkpoint, the system
detects their presence through geofencing and confirms the visit; only then
is the location-specific media for that checkpoint unlocked --- a video
testimonial, a historical image, or a narrative clip --- rendered anchored in
the real world through the device camera using an AR overlay.

Exploration is reinforced by a gamification layer. Users accumulate AR
Points through quiz answers, milestone completions, and side quests; earn
digital milestone badges for completing all milestone stops of a route;
answer five-question multiple-choice knowledge quizzes unlocked after
finishing a route; and compete on a community leaderboard. Points are
redeemable for tangible campus perks through a reward store, and optional
side quests reward detours that require a minimum five-minute on-site stay.

The platform is administered without code. An admin dashboard provides a
waypoint logger for defining route coordinates, waypoint types, and media; a
place management interface for creating, updating, and deleting visiting
places together with their cascaded route data; quiz management; user
management; and aggregate statistics. Because each place records a visitor
count as visits are confirmed, the system generates engagement analytics
organically, with no additional instrumentation.

\begin{figure}[ht]\centering
  \includegraphics[height=0.8\textheight,width=\textwidth,keepaspectratio]{user-flow}
  \caption{End-to-end user journey: from registration to reward redemption.}\label{fig:userflow}
\end{figure}

\section{Scope}
This build delivers a complete vertical slice of the platform as a web
application, proven end-to-end on a single campus route at Orchid College in
Kathmandu. In scope are: route selection and real-time compass navigation;
geofence-verified checkpoint confirmation; AR media playback anchored in the
camera view; the full gamification loop (AR Points, milestones, side quests,
quizzes, badges, leaderboard, and reward redemption); and the complete admin
tooling required to configure routes, media, quizzes, and users.

Explicitly out of scope for this build are native mobile applications (the
product is browser-based only), offline operation (an active network
connection is required for the real-time proximity stream and media
delivery), and multi-language support (the interface is English-only). The
same infrastructure is designed to be pointed at other premises --- heritage
sites such as Pashupatinath or Boudhanath --- but only the Orchid College
route is demonstrated here.

\section{Report Organisation}
The remainder of this report is organised as follows.
\begin{itemize}
  \item Chapter~\ref{ch:problem-statement} states the problem the platform
        addresses and the motivation for solving it now.
  \item Chapter~\ref{ch:objectives} lists the general and specific
        objectives of the project.
  \item Chapter~\ref{ch:market-analysis} summarises the market context,
        competitive landscape, and target user segments.
  \item Chapter~\ref{ch:design-thinking} summarises the design process
        that shaped the product.
  \item Chapter~\ref{ch:system-architecture} describes the overall system
        architecture, from the browser client to the clustered backend.
  \item Chapter~\ref{ch:database-schemas} documents the database schema and
        the relationships between its collections.
  \item Chapter~\ref{ch:gamification} details the gamification
        engine: points, milestones, side quests, quizzes, badges, and
        rewards.
  \item Chapter~\ref{ch:ar-pathfinder} explains the AR pathfinder ---
        geofencing, the compass HUD, and the camera overlay.
  \item Chapter~\ref{ch:tech-stack} justifies the technology stack
        used across all layers.
  \item Chapter~\ref{ch:results} presents the results achieved during the
        hackathon build.
  \item Chapter~\ref{ch:future-work} discusses current
        limitations and future work.
  \item Chapter~\ref{ch:conclusion} concludes the report.
  \item Appendix~\ref{app:api} provides the API reference,
        Appendix~\ref{app:schemas} the field-level database schema
        reference, and Appendix~\ref{app:code} selected code listings.
\end{itemize}
